\chapter{Introduction}
\label{ch:introduction}

\section{Motivation and Problem Statement}
\label{sec:intro:motivation}

The Internet was designed to deliver packets, not experiences. A router forwards a datagram
toward its destination without any notion of \emph{why} the datagram was sent: whether it
carries a frame of an interactive video call that will be worthless if it arrives late, a
byte of a background software update that only cares about aggregate throughput, or a
control message whose loss would be catastrophic. For decades this deliberate ignorance --
the end-to-end principle -- was a feature. Yet as applications have become more diverse and
more demanding, the mismatch between what the network optimizes and what the application
actually values has grown into one of the central inefficiencies of modern networking.

Path-aware Internet architectures such as SCION~\cite{Perrig2017} change the terms of this
problem. In SCION, endpoints discover \emph{multiple} end-to-end paths to a destination and
choose among them, rather than deferring blindly to whatever route BGP happens to install.
This turns path selection from a network-internal decision into an application-facing one:
for the first time, the entity that knows the application's intent -- the endpoint -- also
holds the lever that determines how traffic is carried. The fundamental question this thesis addresses is:

\begin{quote}
\itshape How should an endpoint choose the best path -- or the best \emph{combination} of
paths -- given both the intention of the application and the current state of the network?
\end{quote}

\subsection{The Research Challenge}
\label{sec:intro:challenge}
Answering this question is hard for three compounding reasons:
\begin{enumerate}
    \item \textbf{``Best'' is not a fixed target:} A path optimal for a bulk transfer (high goodput, tolerant of latency) is a poor choice for a video call (which prizes low latency and jitter). Hand-tuned objectives fail across diverse workloads, and training a separate reinforcement learning (RL) policy for every conceivable objective does not scale.
    \item \textbf{The action space is dynamic and unbounded:} The set of candidate paths differs across source--destination pairs and fluctuates over time as links appear, fail, congest, or recover. Standard policies assuming a fixed, enumerated action set cannot cope.
    \item \textbf{State is partially and expensively observable:} Fresh per-path measurements incur active probing overhead. A practical policy must decide well while measuring little.
\end{enumerate}

\subsection{Key Insights and Approach}
\label{sec:intro:approach}
This thesis addresses these challenges through a unified principle inspired by recent transport-layer auto-tuning. The Mortise system~\cite{Shen2026} shows that network mechanisms can be steered by treating application \emph{intent} as an explicit, dynamic input rather than a constant compiled into an algorithm. We transplant this principle across two sequential operational control layers:

\paragraph{1. Motivating Foundation: Intent-Aware Single-Path Selection.} We first investigate how application intent can modulate discrete path choice for general network flows. Building on the DQN path selector of Bourak~et~al.~\cite{Bourak2025}, we introduce a permutation-equivariant \emph{path-scoring} architecture that supports variable candidate path sets, and which turns out to carry most of the quantitative gain in this chapter: it raises goodput by 41\% over a fixed-action agent, at less than half the parameter count. We then condition the agent on an explicit vector of application intent, and establish a \emph{design lemma} governing where that conditioning may enter: \emph{in comparison-based selection, a signal that reaches only the terms shared by every candidate cancels under the argmax and cannot change the chosen path, no matter how the policy is trained}. The algebra is one line; what makes it worth stating is that a plausible and widely-used routing --- conditioning the value stream of a dueling head --- sits exactly on the wrong side of it, and that the consequence is measurable, with such a policy re-ranking roughly three times less over five training seeds with disjoint intervals. Routing the intent into the \emph{per-candidate} comparison instead lets a single policy serve diverse application intents without retraining, and interpolate to intents it never trained on. This study also exposes two limits: the environment offers only one axis on which intents genuinely conflict, and when congestion surges, even the best single path reaches an architectural performance ceiling.

\paragraph{2. Joint Application--Network Optimization over Multipath.} Driven by the limitations of single-path selection, we expand our framework to co-optimize both the application and the network simultaneously. Leveraging SCION's native path multiplicity and Multipath-QUIC~\cite{DeConinck2017} transport, we address interactive real-time video streaming by controlling both \emph{what} the application produces (the video encoder's target bitrate) and \emph{how} the network carries it (the continuous per-path traffic split), via a two-timescale hierarchical reinforcement-learning system: a slow application agent for bitrate adaptation and a fast transport agent for path scheduling. This is the thesis's best-evidenced claim. The pair beats all seven comparison points under all three application intents, including a WebRTC-style deployed stack, at margins three to seven times the spread between training runs of the same method, and \emph{neither decision can be made well alone} --- removing either agent forfeits most of the achievable quality. The permutation-equivariant head is again what makes it possible rather than merely convenient: substituting a fixed-dimension one costs 0.231 QoE, the largest measured effect in this thesis, and drops the system from first in nine of nine evaluation cells to none.

Two further findings correct intuitions we began with. \emph{What governs how much the joint controller earns is not whether the network is stationary}: it is how unevenly capacity is spread across the paths -- whether an even split fits inside every path's residual capacity -- and how strictly the application prices delay. Path churn and congestion bursts matter because they destroy that evenness transiently, not because motion is the relevant variable; a stationary but asymmetric network defeats application-only adaptation just as thoroughly. And the conditioning principle carries into continuous allocation, where one policy taking the intent as an input replaces three separately trained specialists at a third of the training cost, with no penalty detectable at the resolution our experiments achieve --- but it also finds that principle's boundary. The cancellation result is exact only for scorers additively separable in the conditioning; we measure three ways an ordinary deployed architecture violates that premise, and a trained nonlinear encoder consequently routes intent past a cancellation the theory says should block it. The theory says where conditioning must enter to be \emph{guaranteed} to act; the experiment shows a deep encoder arriving there without a guarantee.

\subsection{Guiding Questions}
\label{sec:intro:questions}
Throughout this investigation, four central puzzles guide the reader:
\begin{itemize}
    \item \emph{Can a single RL policy realize a range of application intents without retraining, and where must the intent enter the network for it to change the policy's choice at all?}
    \item \emph{Where are the exact operational boundaries where single-path selection falls short?}
    \item \emph{When is it worth splitting traffic across multiple paths, and does joint application--network co-optimization strictly beat optimizing either layer in isolation?}
    \item \emph{How do these learned controllers perform regarding probing overhead, real-time inference latency, and deployability?}
\end{itemize}

\section{Contributions}
\label{sec:intro:contributions}

This thesis makes the following explicit, verifiable contributions:

\begin{enumerate}
  \item \textbf{An intent-conditioned single-path selection framework for SCION:} We extend the DQN path selector of Bourak~et~al.~\cite{Bourak2025} with a permutation-equivariant path-scoring network and per-path intent conditioning, enabling a single policy to realize a range of application intents without per-objective retraining, and to serve intents it never trained on \emph{zero-shot}: behavior interpolates monotonically between trained objectives (Spearman $\rho = -1.000$), including on source--destination pairs never seen in training (\cref{ch:single_path}).

  \item \textbf{An analysis of where intent conditioning must enter:} We show analytically that a conditioning signal reaching only the terms shared by every candidate cannot change a comparison-based selection, and confirm empirically, over five training seeds with disjoint confidence intervals, that such a policy re-ranks roughly three times less than one whose conditioning reaches the per-candidate comparison (\cref{ch:single_path}).

  \item \textbf{Verifiable probing reduction at near-oracle quality:} We demonstrate that our learned single-path selector falls short of a per-context reward oracle by at most 0.0018 reward on any evaluated intent --- while inspecting one path per decision rather than all thirty --- and outperforms every heuristic baseline on three of four intents, cutting probing cost by roughly $40\times$ against exhaustive bandwidth measurement and probing less than every heuristic in absolute terms. We also bound what that proximity means: the oracle's own re-ranking rate is 0.161, so most decision contexts admit no better answer under a different intent, and near-oracle performance partly reflects a task with little headroom (\cref{ch:single_path}).

  \item \textbf{A hierarchical system for joint application--network co-optimization:} We design and implement a two-timescale hierarchical SAC framework for real-time video over Multipath QUIC, co-optimizing target bitrates and continuous path splits across NS-3 packet simulation and analytical testbeds. It beats all seven comparison points under all three application intents --- including a WebRTC-style deployed stack, which it leads by 0.065, 0.118, and 0.143 QoE against training-seed spreads of 0.018--0.023 --- ranking first in every evaluation cell under two intents and eight of nine under the third, and delivering higher VMAF at lower delivered latency in each case (\cref{ch:multipath}).

  \item \textbf{The permutation-equivariant path head is load-bearing, not a convenience:} Scoring each candidate with shared weights is usually presented as the way to accommodate a variable action set. We show it is what makes the system work at all. Substituting a conventional fixed-dimension head, with the two configurations otherwise field-identical, costs $0.231 \pm 0.029$ QoE, positive in all nine evaluation cells against training-seed spreads of 0.013 and 0.037 --- the largest measured effect in this thesis, and enough that with the flat head the preceding contribution is simply false, the learned pair ranking first in zero of nine cells rather than nine. The same architecture raises single-path goodput by 41\% over a fixed-action agent, at less than half the parameter count, so on both problems the dynamic-input design is cheaper to store as well as better to run (\cref{ch:single_path,ch:multipath}).

  \item \textbf{Intent conditioning extended to continuous allocation, with its premise measured:} We prove the cancellation result's softmax analogue --- for an additively separable scorer, conditioning entering only path-independent terms leaves the traffic split numerically identical, exactly rather than up to re-ranking --- and show that a single conditioned policy replaces three per-intent specialists at a third of the training cost, with no penalty detectable at the resolution these experiments achieve, adapting continuously across the space between them. We state that as the null it is: the training-seed spread is 0.041--0.059 against effects of order 0.02--0.05, and the bar itself is noise-limited, since the specialists proved not to be measurably specialized. We then measure three mechanisms by which the deployed architecture violates the proposition's premise, so that a nonlinear encoder routes conditioning past the cancellation in practice: the result identifies which routings carry a \emph{guarantee}, not which happen to work (\cref{ch:multipath}).

  \item \textbf{Characterization of when joint multipath control pays off:} We show the governing quantity is not network volatility but the relationship between the load an even split places on each path and that path's residual capacity, together with how strictly the application prices delay. A stationary but asymmetric topology defeats application-only adaptation as thoroughly as a churning one, while on a well-balanced topology the measured margin falls below the spread between training runs of the same method --- so the regime a scheduler is benchmarked in decides whether it can be measured at all (\cref{ch:multipath,sec:analysis}).

  \item \textbf{A measurement discipline for learned network control, arrived at by retraction:} Four results in this thesis were reported from a single training run and did not survive being run three times, because the spread over training seeds exceeds the confidence interval over evaluation seeds by an order of magnitude on this problem. Sampling training seeds is necessary and, we show, not sufficient: across six independently trained policies the analytical backend's ranking is uncorrelated with the packet-level backend's ($\rho = -0.086$, the best surrogate checkpoint being the worst on the packet-level body), so a cheap surrogate can be an adequate training environment while carrying no information for selecting among the policies it produces (\cref{ch:multipath}).
\end{enumerate}

\section{Organization}
\label{sec:intro:organization}

The remainder of this thesis is organized as follows:

\Cref{ch:background} reviews technical foundations, including SCION, real-time video transport, deep RL machinery, and goal-conditioned policies.
\Cref{ch:problem} formalizes the decision problems, system models, desired properties, and trust assumptions.

\Cref{ch:single_path} presents our study on intent-conditioned single-path selection, detailing its scoring architecture, the conditioning result and its proof, and the performance limits it exposes.
\Cref{ch:multipath} expands into joint application--network optimization over multipath, detailing the two-timescale SAC architecture and evaluation across dynamic network regimes.

\gap{say here, where the reader forms expectations about what each chapter contains, that the two
studies do not share a substrate. \Cref{ch:single_path} runs on a generated SCION topology with
beaconing, segments, and a path store; \cref{ch:multipath} runs on a hand-specified six-path
configuration (\cref{tab:app:p2topo}) with no SCION control plane at all, using only the path
multiplicity SCION provides. This is a defensible scoping decision --- the multipath question is about
transport and application control, not about path discovery --- but leaving it implicit invites the
reader to assume a continuity of environment that does not exist, and to be surprised in
\cref{sec:mp:impl}. One sentence here, and the corresponding one in \cref{ch:multipath}'s opening,
settle it.}

\Cref{ch:synthesis} unifies the findings, analyzing overhead, scalability, generalization, and deployment considerations (\cref{sec:analysis,sec:discussion}).
\Cref{ch:related} places our work within the broader research landscape, and \Cref{ch:conclusion} concludes.